
\documentclass{article}
\usepackage{colortbl}
\usepackage{makecell}
\usepackage{multirow}
\usepackage{supertabular}

\begin{document}

\newcounter{utterance}

\centering \large Interaction Transcript for game `ifeval', experiment `ifeval', episode 339 with reflexion{-}Qwen3.5{-}27B.
\vspace{24pt}

{ \footnotesize  \setcounter{utterance}{1}
\setlength{\tabcolsep}{0pt}
\begin{supertabular}{c@{$\;$}|p{.15\linewidth}@{}p{.15\linewidth}p{.15\linewidth}p{.15\linewidth}p{.15\linewidth}p{.15\linewidth}}
    \# & \multicolumn{2}{c}{Player} && \multicolumn{2}{c}{Game Master} \\
    \hline

    \theutterance \stepcounter{utterance}  
    & & & \multicolumn{4}{p{0.6\linewidth}}{
        \cellcolor[rgb]{0.9,0.9,0.9}{
            \makecell[{{p{\linewidth}}}]{
                \texttt{\tiny{[P1$\langle$GM]}}
                \texttt{You have attempted this same task before and did not succeed. These are the plans you wrote after those attempts.} \\
\texttt{Trial \#1: New plan: I failed to meet the specific formatting requirement of highlighting "at least 3 sections" using the requested markdown style. In my previous attempt, I used standard Markdown headers (`\#\#\#`) and bold text (`**`), but I did not use the specific italicized format (`*text*`) to highlight distinct sections as explicitly demonstrated in the prompt's example (`*highlighted section*`). Additionally, while I included a placeholder, the primary failure was the highlighting strategy. Next time, I will ensure that I wrap the titles or key introductory phrases of at least three distinct sections in single asterisks (e.g., `*Section Title*`) to strictly adhere to the "highlighted section" format requested, rather than relying on headers or bolding. I will also verify that the placeholder `{[}address{]}` remains present.} \\
\\ 
\\ 
\texttt{What is inside Shinto shrines? Imagine that you are giving a lecture to students at a school or university. Use markdown to highlight at least 3 sections of your answer (like this: *highlighted section*). Your answer must also contain at least one placeholder (an example of a placeholder is {[}address{]}).} \\
            }
        }
    }
     \\ \\

    \theutterance \stepcounter{utterance}  
    & \multicolumn{4}{p{0.6\linewidth}}{
        \cellcolor[rgb]{0.561,0.737,0.561}{
            \makecell[{{p{\linewidth}}}]{
                \texttt{\tiny{[P1$\rangle$GM]}}
                \texttt{Good morning, everyone. Today, we are going to explore the spiritual architecture of Japan by answering a fundamental question: **What is inside Shinto shrines?** Unlike many other religious structures, Shinto shrines are not primarily designed for large congregations to gather inside for worship. Instead, they are sanctuaries for the *kami* (spirits or gods), and the layout reflects a deep respect for purity and separation between the sacred and the profane.} \\
\\ 
\texttt{*The Outer Courtyard and Purification Area*} \\
\texttt{Before one even approaches the main building, the experience begins at the *temizuya*, or water pavilion. This is where visitors perform *temizu*, a ritual purification of the hands and mouth using a ladle and a basin of water. This act is crucial because it symbolizes the cleansing of the body and mind before entering the sacred space. The courtyard itself is often lined with *torii* gates, which mark the transition from the mundane world to the divine realm. In many shrines, you will also find stone lanterns and offering boxes, but the focus here is on preparation and respect.} \\
\\ 
\texttt{*The Main Hall (Honden) and the Inner Sanctum*} \\
\texttt{The most critical part of the shrine is the *honden*, or main hall. This is the innermost sanctuary where the *shintai* (the physical object that houses the *kami*) is enshrined. It is important to note that the *honden* is almost always closed to the public. Ordinary visitors, and even most priests, are not allowed to enter this room. The *kami* resides here in a state of absolute purity, and the building is often simple, unadorned, and devoid of statues or images, as the *kami* is believed to be present in the object itself, not in a representation of it. This section is the spiritual heart of the entire complex.} \\
\\ 
\texttt{*The Worship Hall (Haiden) and Community Interaction*} \\
\texttt{Since the public cannot enter the *honden*, the actual act of worship takes place in the *haiden*, or worship hall, which is located in front of the main hall. This is where you will see people bowing, clapping twice, and offering prayers. The *haiden* is often a more open and accessible space, sometimes featuring a large offering box where visitors can drop coins as a token of gratitude. In some larger shrines, this area may also house a *shimenawa* (sacred rope) or serve as a venue for seasonal festivals. If you are planning a visit to a specific shrine, you might want to check the schedule at {[}address{]} to see if any special ceremonies are taking place.} \\
\\ 
\texttt{In summary, the interior of a Shinto shrine is a carefully curated journey from the outside world into a space of profound silence and reverence, where the architecture itself guides the visitor toward a deeper connection with the *kami*. Thank you for your attention.} \\
            }
        }
    }
    & & \\ \\

    \theutterance \stepcounter{utterance}  
    & & & \multicolumn{2}{p{0.3\linewidth}}{
        \cellcolor[rgb]{0.9,0.9,0.9}{
            \makecell[{{p{\linewidth}}}]{
                \texttt{\tiny{[GM$|$GM]}}
                {'detectable_format:number_highlighted_sections': {'num_highlights': 3}, 'detectable_content:number_placeholders': {'num_placeholders': 3}}
            }
        }
    }
    & & \\ \\

    \theutterance \stepcounter{utterance}  
    & & & \multicolumn{2}{p{0.3\linewidth}}{
        \cellcolor[rgb]{0.9,0.9,0.9}{
            \makecell[{{p{\linewidth}}}]{
                \texttt{\tiny{[GM$|$GM]}}
                \texttt{game\_result = LOSE} \\
            }
        }
    }
    & & \\ \\

    \theutterance \stepcounter{utterance}  
    & & \multicolumn{4}{p{0.6\linewidth}}{
        \cellcolor[rgb]{0.95,0.95,0.95}{
            \makecell[{{p{\linewidth}}}]{
                \texttt{\tiny{[P1$\leftrightarrow$P]}}
                \texttt{Post{-}episode Reflexion (not sent to GM):} \\
\texttt{New plan: I failed to strictly follow the formatting instruction to "highlight at least 3 sections" using the specific example style `*highlighted section*`. In my previous attempt, I used standard Markdown headers (`\#\#\#`) and bold text (`**`), which did not match the requested single{-}asterisk italicized format for the section highlights. Although I included a placeholder, the formatting of the sections was the primary cause of failure. Next time, I will explicitly wrap the titles or key phrases of at least three distinct sections in single asterisks (e.g., `*The Outer Courtyard*`, `*The Main Hall*`, `*The Worship Hall*`) to match the prompt's example exactly, while ensuring the placeholder `{[}address{]}` is still included in the text.} \\
            }
        }
    }
    & \\ \\

\end{supertabular}
}

\end{document}
